\chapter{Lista 13.}

\begin{problem}[1]
		Udowodnić, że pierścień $K\left[ X^2, X^3 \right]$ nie jest UFD.     
\end{problem} 
\begin{solution}
    TODO
\end{solution} 
\begin{problem}[2]
    Niech $R$ będzie UFD i $a,b \in R$. Udowodnić, że ideał $\left( a \right) \cap \left( b \right)$ jest główny.   
\end{problem} 
\begin{solution}
    TODO
\end{solution} 
\begin{problem}[3]
    Niech $S$ będzie podzbiorem multiplikatywnym $R$. Udowodnić, że:
\begin{enumerate}[label=(\alph*)]
    \item dla $\mathcal{I} \trianglelefteqslant R_{S} $ zachodzi $\left( \mathcal{I}\cap R \right) R_{S} = \mathcal{I}$.
		\item istnieje naturalny monomorfizm $f: R_{S}\to R_{0}$.
		\item jeśli $R$ jest PID, to $R_{S}$ też jest PID.        
\end{enumerate}
\end{problem} 
\begin{solution}
    TODO
\end{solution} 
\begin{problem}[4]
    Wywnioskować z poprzedniego zadania, że $\mathbb{Z}\left[ \tfrac{1}{2}\right]$ jest UFD. 
\end{problem} 
\begin{solution}
    TODO
\end{solution} 
\begin{problem}[5]
    Dla każdej liczby pierwszej $p$ w naturalny sposób utożsamiamy $\mathbb{Z}_{(p)}$ z podpierścieniem ciała $\mathbb{Q}$. Udowodnić, że 
\[
    \bigcap_{p \in \mathbb{P}} \mathbb{Z}_{(p)} = \mathbb{Z}    
,\]   
    gdzie $\mathbb{P}$ jest zbiorem wszystkich liczb pierwszych. 
\end{problem} 
